\documentclass[10pt,a4paper]{amsart}
\usepackage[margin=19mm]{geometry}
\usepackage{amsmath,amssymb,mathtools}
\usepackage[scr=rsfs,cal=euler]{mathalfa}
\usepackage{xfrac}
\usepackage[unicode,hidelinks,bookmarks=false]{hyperref}
\newcommand{\ie}{i.e.\ }
\newcommand{\cf}{cf.\ }
\newcommand{\resp}{respectively\ }
\newcommand{\Subsection}{\subsection}
\providecommand{\nobreakdash}{\nobreak\hbox{-}}
\setlength{\emergencystretch}{2em}
\multlinegap0pt
\usepackage{bm}
\usepackage{amssymb}
\usepackage{enumerate}
\usepackage{float}
\newtheorem{theorem}{Theorem}
\newtheorem{lemma}{Lemma}
\newtheorem{proposition}{Proposition}
\newtheorem{corollary}{Corollary}
\def\E{\mathbb{E}}
\def\T{\mathbb{T}_L}
\def\d{\partial}
\def\f{\varphi}
\def\tu{\tilde{u}}
\title{Long Time Behavior of Stochastic Thin Film Equation}
\author{Oleksiy Kapustyan, Olha Martynyuk, Oleksandr Misiats, Oleksandr Stanzhytskyi}
\date{Source: arXiv:2604.10010v1 (CC BY 4.0)}
\begin{document}
\maketitle

\noindent\emph{Source and licence.} Long Time Behavior of Stochastic Thin Film Equation. Version arXiv:2604.10010v1 (CC BY 4.0), distributed by arXiv under the Creative Commons Attribution 4.0 International licence, http://creativecommons.org/licenses/by/4.0/, as recorded in the arXiv licence metadata for this exact version and shown on the abstract page https://arxiv.org/abs/2604.10010v1. Full licence notice: \texttt{licenses/arxiv-2604.10010-CC-BY-4.0.txt}.\par\smallskip

\noindent\textit{Editorial scope.} Counted unit: the article's proposition prop:5.1 (source0.tex lines 1603--1612). The article's complete original proof of it is delivered with it (source0.tex lines 1614--1800, one \texttt{proof} environment of 187 lines). No mathematics is written, repaired or re-derived: the statement and the proof are the article's own lines, in the article's own notation, and no retained line is substituted or re-flowed.  Retained uncounted as the source-local context the proof needs: lines 223--230; lines 403--411; lines 415--427; lines 670--710; lines 1192--1223; lines 1251--1262; lines 1403--1411; lines 1415--1419.  Nothing was omitted from any retained span, so the package carries no omission record.  No retained line contains a citation, so the wrapper carries no bibliography and no bibliography entry.  No retained line contains a figure environment, an image-inclusion command or drawing code, so no image is delivered and the package declares no asset dependency.  Editorial formatting only, in addition: an AMS \texttt{amsart} preamble that restates the article's own declarations and macros used by the retained lines, namely the environments \texttt{theorem}, \texttt{lemma}, \texttt{proposition}, \texttt{corollary} on separate counters, together with the standard packages those lines need.  The retained environments print in this document as Theorem 1, Lemma 1, Proposition 1, Corollary 1 in order of appearance, whereas the article prints the same items with the numbering of its own full document.  The complete native source of the article as distributed in the arXiv e-print for this exact version is bundled unchanged as \texttt{sources/198242/source0.tex}.
\begin{equation}
\mathrm{d}u
=
- \partial_x\!\left(u^n \partial_x^3 u\right)\,\mathrm{d}t
+ \gamma(t) u \,\mathrm{d}t
+ \alpha(t) u \,\mathrm{d} \beta (t),
\label{1.1}
\end{equation}

\begin{theorem}\label{Th:2.1}
Assume $n > 0$. Then for any
$u_0 \in H^1(\mathbb{T})$, $u_0 \ge 0$, $u_0 \not\equiv 0$,
equation \eqref{1.1} admits a nonnegative weak martingale solution $\tilde{u}$, which is defined on $t \in [0,T]$ for any $T > 0$, and for any $p \geq 2$ satisfies
\[
\tilde{\E} \sup_{t \in [0,T]} \|\tilde{u}(t,\cdot)\|_2^p \leq C \| u_0\|^p_2
\]
for some $C>0$ independent of $\tilde{u}$ and $u_0$.
\end{theorem}

\begin{theorem}\label{th:2.2}
Assume $n>2$ and $u_0 \in H^1(\mathbb{T}_L)$, is such that $u_0 \ge 0$, and
\[
\int_L u_0^{2-n}(x) \, dx < \infty.
\]
Then the equation~\eqref{1.1} admits a nonnegative martingale solution
$\tu(t)$ on $[0,\infty)$ such that for any $t>0$,
the set
\[
\{x \in L : \tu(t,x,\omega)=0\}
\]
has zero Lebesgue measure, $\tilde{\mathbb{P}}$--almost surely.
\end{theorem}

\begin{lemma}
\label{lem:3.4}
Let $p \in [2,\infty)$ and
\[
w(t_0,x): = w_0 \in L^p(\Omega,\mathcal{F}_0,\mathbb{P}; H^1(\mathbb{T}_L)).
\]
Then for all $0 \leq t_0 \leq t \leq T$ the following properties hold:
\begin{enumerate}
\item
\begin{equation}
\label{3.6}
\mathbb{E} \sup_{s \in [t_0,t]} \|w(s, \cdot)\|_{1,2}^p
\le e^{C_1(t-t_0)} \mathbb{E} \|w (t_0, \cdot)\|_{1,2}^p.
\end{equation}

\item
If $w_0 \ge 0$ almost surely, then
$w(t,x) \ge 0$ almost surely for all $t \in [0,T]$.



\item 
\begin{equation}
\label{3.7}
\mathbb{E}\|\d_x w(t, \cdot)\|_{L^2}^p
\le
\mathbb{E}\|\d_x w(t_0, \cdot)\|_{L^2}^p \,
e^{C_2 (t-t_0)}.
\end{equation}

\item
\begin{equation}
\label{3.7.1}
\mathbb{E}\left|\int_{L} w(t,x)\,dx\right|^p
\le
\mathbb{E}\left|\int_{L} w(t_0,x)\,dx\right|^p
e^{C_1 (t-t_0)}.
\end{equation}
\end{enumerate}
Here $C_1>0$ is a constant independent of $w$.
\end{lemma}

\begin{proposition}\label{prop:3.7}
Denote
\[
X_u := BC^0([0,T]\times \T), \quad
X_J := L^2([0,T]\times \T)\ \text{(with weak topology)}, \quad
X_\beta := BC^0([0,T];\mathbb{R}).
\]
Then there exist random variables
\[
\tilde u, \tilde u_N : [0,T] \to X_u, \qquad
\tilde J, \tilde J_N : [0,T] \to X_J, \qquad
\tilde \beta_N^{'}, \tilde \beta_N : [0,T] \to X_\beta,
\]
with
\[
(\tilde u_N, \tilde J_N, \tilde \beta_N^{'}) \sim (u_N, J_N, \beta_N),
\]
where
\[
J_N := \chi_{v_N > 0}\, v_N^n(\partial_x^3 v_N).
\]
Furthermore, there exists a subsequence (still indexed by $N$) such that
\[
\tilde u_N(\omega) \to \tilde u(\omega) \quad \text{in } X_u, \qquad
\tilde J_N(\omega) \rightharpoonup \tilde J(\omega) \quad \text{in } X_J,
\]
and
\[
\tilde \beta_N^{'}(\omega) \to \tilde \beta(\omega) \quad \text{in } X_\beta,
\]
for every $\omega \in [0,1]$, as $N \to \infty$.
\end{proposition}

\begin{corollary}\label{cor:3.8}
For $\tilde u_N$, $\tilde J_N$, $\tilde W_N$ and $\tilde u$ as in Proposition \ref{prop:3.7}, we have
\begin{equation}\label{3.19}
\|\tilde u_N - \tilde u\|_{BC^0([0,T]\times \T)} \to 0,
\quad
\|\tilde v_N - \tilde u\|_{L^\infty([0,T]\times \T)} \to 0,
\quad
\|\tilde w_N - \tilde u\|_{L^\infty([0,T]\times \T)} \to 0,
\tag{3.19}
\end{equation}
as $N \to \infty$, $\tilde{\mathbb{P}}$-almost surely.
\end{corollary}

\begin{lemma}\label{lem:4.1}
For any $t \in [0,T]$ we have
\begin{equation}
\mathbb{E} \int_L G(v_n(t,x))\, dx 
\le C \int_\Omega G(u_0(x))\, dx,
\label{4.1}
\end{equation}
for some constant $C>0$ independent of $N$ and $u_0$.
\end{lemma}

\begin{equation}
\int_L G(v_n(t,x))\, dx 
\le \int_L G(v_n((j-1)\delta,x))\, dx.
\label{4.2}
\end{equation}

\begin{proposition}\label{prop:5.1}
Suppose $n \ge 4$ and $u_0 \in H^{1}(\mathbb{T}_L)$ with
$u_0(x) > 0$ on $[0,1]$. Then there exists a set
$A \subset \Omega$ with $\mathbb{P}(A)=1$ such that
for all $\omega \in A$,
\[
u(t,x,\omega) > 0
\quad \text{for all } t \ge 0, \; x \in L.
\]
\end{proposition}

\begin{proof}
In order to establish this proposition, we need to improve the Lemma \ref{lem:4.1} estimate on $u^{n}(t,x,\omega)$. Fix a an arbitrary interval $[0,T]$. According to Theorem \ref{th:2.2}, the equation \eqref{1.1} admits a weak martingale solution on $[0,\infty)$. Moreover, this solution is non-negative, and for every $t>0$ the set
\[
\{ x \in L : u(t,x,\omega)=0 \}
\]
has zero Lebesgue measure $\mathbb{P}$ - almost surely.
For $n \ge 4$, consider the approximating sequences
$v_N$ and $w_n$, constructed in Theorem \ref{Th:2.1}.

For any $j=1,\dots,N+1$ and
$t \in [(j-1)\delta,j\delta)$, by \eqref{4.2} we have 
\begin{equation*}
\int_{L} v_N^{2-n}(t,x)\,dx
\le
\int_{L} v_N^{2-n}((j-1)\delta,x)\,dx.
\label{5.1}
\end{equation*}
Furthermore, applying It\^o's formula,
\begin{align*}
\int_{L} w_N^{2-n}(t,x)\,dx
&=
\int_{L} w_N^{2-n}((j-1)\delta+,x)\,dx \notag \\
&\quad
+ \int_{(j-1)\delta}^{t}
\Big[(2-n)\gamma(s)
+ \frac{(1-n)(2-n)}{2}\alpha^2(s)\Big]
\int_{L} w_N^{2-n}(s,x)\,dx\,ds \notag \\
&\quad
+ \int_{(j-1)\delta}^{t}
(2-n)\alpha(s)
\int_{L} w_N^{2-n}(s,x)\,dx\, d\beta(s).
\label{5.2}
\end{align*}
Therefore, $\displaystyle \int_{L} w_n^{2-n}(t,x)\,dx$
satisfies a linear stochastic differential equation.
Hence,
\begin{align*}
&\int_{L} w_N^{2-n}(t,x)\,dx \\
& =
\int_{L} w_N^{2-n}((j-1)\delta+,x)\,dx  \notag 
\quad \times
\exp\Bigg\{
\int_{(j-1)\delta}^{t}
\varphi(s)\,ds
+
\int_{(j-1)\delta}^{t}
\psi(s)\,d\beta(s)
\Bigg\},
\end{align*}
where
\[
\f(s): = (2-n)\gamma(s)
+ \frac{(1-n)(2-n)}{2}\alpha^2(s) - \frac{(2-n)\alpha(s)}{2},
\]
and 
\[
\psi(s):= (2-n) \alpha(s).
\]
Consequently, using the argument of Lemma \ref{lem:4.1}, on $[0, \delta)$ we obtain
\begin{equation*}
\int_{L} u_N^{2-n}(t,x)\,dx
\le
\int_{L} u_0^{2-n}(x)\,dx.
\end{equation*}
Hence,
\begin{align*}
    \int_{L} w_N^{2-n}(t,x)\,dx
& = \int_{L} w_N^{2-n}(0,x)\,dx
\exp\left( \int_0^t \varphi(s)\,ds + \int_0^t \psi(s)\,d\beta(s) \right) \\
& = \int_{L} v_N^{2-n}(\delta-,x)\,dx
\exp\left( \int_0^t \varphi(s)\,ds + \int_0^t \psi(s)\,d\beta(s) \right) \\
& \leq \int_{L} u_0^{2-n}(x)\,dx
\exp\left( \int_0^t \varphi(s)\,ds + \int_0^t \psi(s)\,d\beta(s) \right).
\end{align*}
In a similar way, on $[\delta,2\delta)$ we obtain
\begin{align*}
& \int_{L} v_N^{2-n}(t,x)\,dx \le \int_{L} v_N^{2-n}(\delta+,x)\,dx = \int_{L} w_N^{2-n}(\delta-,x)\,dx  \\
& = \int_{L} w_N^{2-n}(0,x)\,dx
\exp\left( \int_0^\delta \varphi(s)\,ds + \int_0^\delta \psi(s)\,d\beta(s) \right)\\ 
& =  \int_{L} v_N^{2-n}(\delta-,x)\,dx
\exp\left( \int_0^\delta \varphi(s)\,ds + \int_0^\delta \psi(s)\,d\beta(s) \right)\\
& \le \int_{L} u_0^{2-n}(x)\,dx
\exp\left( \int_0^\delta \varphi(s)\,ds + \int_0^\delta \psi(s)\,d\beta(s) \right) \\
& \le \int_{L} u_0^{2-n}(x)\,dx
\exp\left( \int_0^t \varphi(s)\,ds + \int_0^t \psi(s)\,d\beta(s) \right).
\end{align*}
On $[2\delta,3\delta]$ we get
\begin{align*}
& \int_{L} v_N^{2-n}(t,x)\,dx \le \int_{L} v_N^{2-n}(2\delta+,x)\,dx = \int_{L} w_N^{2-n}(2\delta-,x)\,dx  \\
& = \int_{L} w_N^{2-n}(\delta+,x)\,dx
\exp\left( \int_\delta^t \varphi(s)\,ds + \int_\delta^t \psi(s)\,d\beta(s) \right)\\ 
& =  \int_{L} u_N^{2-n}(2\delta-,x)\,dx
\exp\left( \int_\delta^t \varphi(s)\,ds + \int_\delta^t \psi(s)\,d\beta(s) \right)\\
& \le \int_{L} u_0^{2-n}(x)\,dx
\exp\left( \int_0^t \varphi(s)\,ds + \int_0^t \psi(s)\,d\beta(s) \right).
\end{align*}
Therefore, for any $t \in [(j-1)\delta, j\delta)$,
\[
\int_{L} v_N^{2-n}(t,x)\,dx
\le
\int_{L} u_0^{2-n}(x)\,dx
\exp\left( \int_0^t \varphi(s)\,ds + \int_0^t \psi(s)\,d\beta(s) \right).
\]
Hence,
\[
\sup_{t \in [0,T]}
\int_{L} v_N^{2-n}(t,x)\,dx
\le
\int_{L} u_0^{2-n}(x)\,dx
\sup_{t \in [0,T]}\left[
\exp\left( \int_0^t \varphi(s)\,ds + \int_0^t \psi(s)\,d\beta(s) \right) \right].
\]
Next, using the same argument as in the proof of Lemma \ref{lem:3.4}, the random process
\[
\xi(t)
:=
\exp\left(
\int_0^t \varphi(s)\,ds
+
\int_0^t \psi(s)\,d\beta(s)
\right)
\]
is a solution of a linear equation
\[
\xi(t)
= 1
+ \int_{0}^{t}
\Big[
(2-n) \gamma(s)
+ \frac{(2-n)(1-n)}{2}\,\alpha^2(s)
\Big] \xi(s)\,ds
+ \int_{0}^{t} (2-n)\, \alpha(s) \,\xi(s)\,d\beta(s).
\]
This way,
\[
\mathbb{E}\,\sup_{t\in[0,T]} \xi(t)
\le
\Big(\mathbb{E}\,\sup_{t\in[0,T]} \xi^2(t)\Big)^{1/2}
\le C(T),
\]
where we used the martingale properties of stochastic integrals.
Consequently,
\[
\mathbb{E}\,\sup_{t\in[0,T]}
\int_{L} v_N^{2-n}(t,x)\,dx
\le
\int_{L} u_0^{2-n}(x)\,dx \; C(T).
\]
Using Fatou's lemma, Corollary \ref{cor:3.8}, and Theorem \ref{th:2.2}, we obtain
\begin{equation}\label{5.5}
\mathbb{E}\,\sup_{t\in[0,T]}
\int_{L} u_N^{2-n}(t,x)\,dx
\le
\int_{L} u_0^{2-n}(x)\,dx \; C(T).
\end{equation}
In order to complete the proof of Proposition \ref{prop:5.1}, we now argue by contradiction. Assume there exists a set $B \subset \Omega$ such that $P(B)>0$ and
for any $\omega\in B$ there exist $t(\omega)$, $x_0(\omega)$ such that
\[
u(t(\omega),x_0(\omega),\omega)=0.
\]
Using the embedding $H^1(\T) \subset C^{1/2}(\T)$, % (?)
we have
\[
|u(t(\omega),x,\omega) - u(t(\omega),x_0(\omega),\omega)|
\le
K(t,\omega)\,|x-x_0(\omega)|^{1/2}.
\]
Consequently, 
\begin{align}\label{5.6}
 \nonumber &  \mathbb{E}\,\sup_{t\in[0,T]}
\int_{K} u^{2-n}(t,x)\,dx
\ge
\int_{\{\omega\in B\}}
\sup_{t\in[0,T]}
\int_{L} u^{2-n}(t,x,\omega)\,dx \; P(d\omega) \\
& \ge \int_{\{\omega\in B\}}
\int_{L} u^{2-n}(t(\omega),x,\omega)\,dx \; P(d\omega) \\
\nonumber & \ge
\int_{\{\omega\in B\}}
\frac{1}{K^{\,2-n}(t,\omega)}
\int_{L}
\frac{dx}{|x-x_0(\omega)|^{\frac{n-2}{2}}}
\,P(d\omega)
= \infty.
\end{align}
The latter integral is divergent since $n \ge 4$. Thus, \eqref{5.6} contradicts   \eqref{5.5}, and the proof of Proposition \ref{prop:5.1} is complete.
\end{proof}

\end{document}
